\documentclass[11pt]{article}

\usepackage[T1]{fontenc}
\usepackage[spanish,es-nodecimaldot]{babel}
\usepackage{amsmath,amssymb,mathtools,bm}
\usepackage{booktabs}
\usepackage{geometry}
\usepackage{enumitem}
\usepackage[hidelinks]{hyperref}

\geometry{margin=1in}

\newcommand{\R}{\mathbb{R}}
\newcommand{\argmin}{\operatorname*{arg\,min}}
\newcommand{\diag}{\operatorname{diag}}

\title{Metodos numericos para localizar minimos de smoothness\\
en la perdida de forecast}
\author{Nota tecnica del proyecto \texttt{paper\_numerical-smoothness-selection}}
\date{}

\begin{document}
\maketitle

\begin{abstract}
Esta nota desarrolla la parte numerica que queda escondida detras de
encontrar los minimos locales de la perdida en smoothness. El procedimiento
implementado no es una simple busqueda sobre una grilla. La tendencia
penalizada y el forecast tienen primeras y segundas derivadas analiticas con
respecto a \(\lambda\); estas derivadas se transforman a la coordenada
normalizada \(S\), se usa una busqueda adaptativa para detectar intervalos que
pueden contener puntos estacionarios y, finalmente, cada raiz de la derivada
se refina con Brent. Tambien se tratan explicitamente los endpoints
\(S=0\) y \(S=1\), se deduplican soluciones numericamente equivalentes y se
clasifican los puntos por el signo de la segunda derivada.

El objetivo es documentar exactamente la matematica y las decisiones numericas
de las rutinas activas del repositorio.
\end{abstract}

\tableofcontents

\section{Problema numerico}

Para una configuracion fija
\[
(d,L,h),
\]
consideramos una ventana de entrenamiento
\[
\bm x=(x_1,\ldots,x_L)^\top
\]
y un bloque futuro observado
\[
\bm z=(z_1,\ldots,z_h)^\top.
\]

La tendencia penalizada es
\begin{equation}
\widehat{\bm\tau}_\lambda
=
\argmin_{\bm\tau\in\R^L}
\left\{
\|\bm x-\bm\tau\|_2^2
+
\lambda\|D_d\bm\tau\|_2^2
\right\}.
\label{eq:pls}
\end{equation}

Sea
\[
Q_d=D_d^\top D_d.
\]
Entonces
\begin{equation}
\widehat{\bm\tau}_\lambda
=
H_\lambda\bm x,
\qquad
H_\lambda=(I+\lambda Q_d)^{-1}.
\label{eq:smoother}
\end{equation}

Despues de continuar la tendencia \(h\) pasos, obtenemos
\[
\widehat{\bm z}_\lambda.
\]
La funcion objetivo de una Validation 1 es
\begin{equation}
F(\lambda)
=
\frac{1}{h}
\|\bm z-\widehat{\bm z}_\lambda\|_2^2.
\label{eq:forecast-loss}
\end{equation}

La funcion que realmente exploramos es
\[
\widetilde F(S)=F(\lambda(S)),
\qquad
0\le S\le1.
\]
La tarea numerica es recuperar todos los minimos locales relevantes
\[
\mathcal M
=
\left\{
S:
\widetilde F_S(S)=0,
\quad
\widetilde F_{SS}(S)>0
\right\},
\]
mas los posibles minimos de frontera.

\section{Representacion espectral del suavizador}

Sea
\[
Q_d
=
U\diag(\delta_1,\ldots,\delta_L)U^\top,
\qquad
\delta_j\ge0.
\]
Para una penalizacion por diferencias de orden \(d\), exactamente \(d\)
eigenvalores son cero.

Definimos
\[
\bm c=U^\top\bm x
\]
y
\begin{equation}
\alpha_j(\lambda)
=
\frac{1}{1+\lambda\delta_j}.
\label{eq:alpha}
\end{equation}
Entonces
\begin{equation}
\widehat{\bm\tau}_\lambda
=
U\diag(\alpha_j(\lambda))\bm c.
\label{eq:spectral-trend}
\end{equation}

La dependencia en \(\lambda\) queda separada en los escalares
\(\alpha_j\). Esto permite obtener derivadas exactas sin diferencias finitas.

\subsection{Primera y segunda derivada}

De \eqref{eq:alpha},
\[
\alpha_j'(\lambda)
=
-
\frac{\delta_j}{(1+\lambda\delta_j)^2}
=
-\delta_j\alpha_j^2,
\]
y
\[
\alpha_j''(\lambda)
=
\frac{2\delta_j^2}{(1+\lambda\delta_j)^3}
=
2\delta_j^2\alpha_j^3.
\]
Por tanto,
\begin{align}
\widehat{\bm\tau}_\lambda'
&=
U\diag(-\delta_j\alpha_j^2)\bm c,
\label{eq:trend-first}
\\
\widehat{\bm\tau}_\lambda''
&=
U\diag(2\delta_j^2\alpha_j^3)\bm c.
\label{eq:trend-second}
\end{align}

Equivalentemente,
\begin{align}
\widehat{\bm\tau}_\lambda'
&=
-H_\lambda Q_d\widehat{\bm\tau}_\lambda,
\\
\widehat{\bm\tau}_\lambda''
&=
2H_\lambda Q_dH_\lambda Q_d
\widehat{\bm\tau}_\lambda.
\end{align}

Estas son las identidades implementadas en
\texttt{pure\_trend\_derivatives}.

\section{Derivadas del forecast}

La continuacion de diferencias finitas puede escribirse como un operador lineal
\[
C_{d,h,L}\in\R^{h\times L}
\]
tal que
\begin{equation}
\widehat{\bm z}_\lambda
=
C_{d,h,L}
\widehat{\bm\tau}_\lambda.
\label{eq:forecast-operator}
\end{equation}

Como \(C_{d,h,L}\) no depende de \(\lambda\),
\begin{align}
\widehat{\bm z}_\lambda'
&=
C_{d,h,L}\widehat{\bm\tau}_\lambda',
\\
\widehat{\bm z}_\lambda''
&=
C_{d,h,L}\widehat{\bm\tau}_\lambda''.
\end{align}

\section{Derivadas exactas de la MSE}

Sea
\[
\bm r_\lambda
=
\bm z-\widehat{\bm z}_\lambda.
\]
La perdida es
\[
F(\lambda)
=
\frac{1}{h}
\bm r_\lambda^\top\bm r_\lambda.
\]

Como
\[
\bm r_\lambda'=-\widehat{\bm z}_\lambda',
\]
se obtiene
\begin{equation}
F'(\lambda)
=
-
\frac{2}{h}
\bm r_\lambda^\top
\widehat{\bm z}_\lambda',
\label{eq:mse-first}
\end{equation}
y
\begin{equation}
F''(\lambda)
=
\frac{2}{h}
\left[
(\widehat{\bm z}_\lambda')^\top
\widehat{\bm z}_\lambda'
-
\bm r_\lambda^\top
\widehat{\bm z}_\lambda''
\right].
\label{eq:mse-second}
\end{equation}

Asi, cada evaluacion produce
\[
\boxed{
F(\lambda),
\qquad
F'(\lambda),
\qquad
F''(\lambda)
}.
\]

No se aproximan estas derivadas mediante diferencias finitas. Esto evita
escoger un paso numerico auxiliar y reduce el numero de forecasts necesarios.

\section{Objetivo rolling y vectorizacion}

Con varios rolling origins, sea \(\mathcal R\) su conjunto y
\(\bm r_{r,\lambda}\) el residual del origin \(r\). La implementacion usa
\begin{equation}
F_{\mathrm{roll}}(\lambda)
=
\frac{
\sum_{r\in\mathcal R}
\|\bm r_{r,\lambda}\|_2^2
}{
N_{\mathrm{score}}
},
\qquad
N_{\mathrm{score}}
=
\sum_{r\in\mathcal R}h_r.
\label{eq:rolling-loss}
\end{equation}

Las derivadas se agregan de la misma forma. Para ventanas y horizontes fijos,
el codigo prepara una sola vez el espectro de \(Q_d\), las ventanas de training
en la base espectral, el operador de forecast y todos los targets futuros.

Si \(B\) contiene las ventanas historicas en coordenadas espectrales y
\[
G=C_{d,h,L}U,
\]
entonces las predicciones de todos los origins pueden escribirse
simultaneamente como
\begin{equation}
\widehat Z_\lambda
=
\left(
B\odot\bm\alpha(\lambda)^\top
\right)G^\top.
\label{eq:vectorized}
\end{equation}
Para las derivadas se reemplaza \(\bm\alpha\) por
\(\bm\alpha'\) o \(\bm\alpha''\).

\section{La coordenada normalizada de smoothness}

Para \(d>0\),
\begin{equation}
S(\lambda)
=
1-
\frac{1}{L-d}
\sum_{j:\delta_j>0}
\frac{1}{1+\lambda\delta_j}.
\label{eq:smoothness}
\end{equation}

Sus primeras dos derivadas son
\begin{equation}
S'(\lambda)
=
\frac{1}{L-d}
\sum_{j:\delta_j>0}
\frac{\delta_j}
{(1+\lambda\delta_j)^2}
>0,
\label{eq:s-first}
\end{equation}
y
\begin{equation}
S''(\lambda)
=
-
\frac{2}{L-d}
\sum_{j:\delta_j>0}
\frac{\delta_j^2}
{(1+\lambda\delta_j)^3}
<0.
\label{eq:s-second}
\end{equation}

Por tanto,
\[
S:[0,\infty]\to[0,1]
\]
es monotona creciente con
\[
S(0)=0,
\qquad
S(\infty)=1.
\]

\section{Regla de la cadena de \(\lambda\) a \(S\)}

Definimos
\[
\widetilde F(S)=F(\lambda(S)).
\]
Como \(S'(\lambda)>0\),
\[
\frac{d\lambda}{dS}
=
\frac{1}{S'(\lambda)}.
\]
Entonces
\begin{equation}
\widetilde F_S(S)
=
\frac{F'(\lambda)}
{S'(\lambda)},
\label{eq:chain-first}
\end{equation}
y
\begin{equation}
\widetilde F_{SS}(S)
=
\frac{F''(\lambda)}
{[S'(\lambda)]^2}
-
\frac{
F'(\lambda)S''(\lambda)
}{
[S'(\lambda)]^3
}.
\label{eq:chain-second}
\end{equation}

Como \(S'(\lambda)>0\),
\begin{equation}
\widetilde F_S(S)=0
\quad\Longleftrightarrow\quad
F'(\lambda)=0.
\label{eq:stationary-equivalence}
\end{equation}

Ademas, en un punto estacionario,
\[
F'(\lambda)=0,
\]
y por tanto
\begin{equation}
\widetilde F_{SS}(S)
=
\frac{F''(\lambda)}
{[S'(\lambda)]^2}.
\label{eq:curvature-equivalence}
\end{equation}
El signo de la curvatura se conserva. Un minimo en \(\lambda\) sigue siendo un
minimo en \(S\).

\section{Como se calcula \(\lambda(S)\)}

Para \(0<S<1\), la implementacion invierte \eqref{eq:smoothness} por
biseccion monotona. Se resuelve
\begin{equation}
\phi(\lambda)
=
S(\lambda)-S_0
=
0.
\label{eq:inverse-problem}
\end{equation}

El algoritmo actual es:

\begin{enumerate}[leftmargin=2.4em]
    \item Si \(S_0=0\), regresar \(\lambda=0\).
    \item Si \(S_0=1\), regresar \(\lambda=\infty\).
    \item Iniciar
    \[
    \lambda_{\mathrm{lo}}=0,
    \qquad
    \lambda_{\mathrm{hi}}=1.
    \]
    \item Mientras
    \[
    S(\lambda_{\mathrm{hi}})<S_0,
    \]
    usar
    \[
    \lambda_{\mathrm{hi}}
    \leftarrow
    10\lambda_{\mathrm{hi}},
    \]
    hasta encerrar el target o alcanzar \(10^{16}\).
    \item Bisecar el intervalo y conservar la mitad que contiene el target.
    \item Detener cuando
    \[
    |S(\lambda_{\mathrm{mid}})-S_0|<10^{-11}
    \]
    o despues de \(100\) iteraciones.
\end{enumerate}

La biseccion explota una propiedad fuerte del problema: la monotonicidad de
\(S(\lambda)\) esta demostrada, por lo que la raiz interior es unica.

\section{Por que no basta con un optimizador local}

Si \(\widetilde F(S)\) fuera unimodal, un metodo local podria ser suficiente.
Pero aqui queremos recuperar
\[
S_1^\star,
S_2^\star,
\ldots,
S_J^\star,
\]
no solamente
\[
\argmin_S \widetilde F(S).
\]

Un unico Newton, BFGS o golden-section search puede encontrar un minimo y
omitir las demas cuencas. Por ello la implementacion separa
\[
\text{descubrimiento de regiones con raices}
\quad\longrightarrow\quad
\text{refinamiento de cada raiz}.
\]

\section{Dominio interior de busqueda}

La busqueda derivativa se hace sobre
\begin{equation}
0\le S\le1-\eta,
\qquad
\eta=10^{-6}.
\label{eq:interior-domain}
\end{equation}

El punto \(S=1\) corresponde a \(\lambda=\infty\) y se trata por separado
como endpoint.

\section{Malla inicial y refinamiento de endpoints}

La rutina inicia con
\[
m_0=9
\]
puntos uniformes:
\[
S_i
=
\frac{i}{m_0-1}(1-\eta),
\qquad
i=0,\ldots,m_0-1.
\]

No es una grilla final. Solo es el esqueleto de una busqueda adaptativa.

La compactificacion puede comprimir estructura cerca de \(S=1\), y en menor
grado cerca de \(S=0\). Por eso se insertan seis niveles diadicos dentro de las
dos celdas de frontera. Si \(w\) es el ancho de una celda inicial, se agregan
distancias
\[
\frac{w}{2},
\frac{w}{2^2},
\ldots,
\frac{w}{2^6}
\]
desde cada extremo.

\section{Refinamiento adaptativo}

Definimos
\[
g(S)=\widetilde F_S(S),
\qquad
q(S)=\widetilde F_{SS}(S).
\]
Para un intervalo \([a,b]\), sea
\[
m=\frac{a+b}{2}.
\]
Evaluamos
\[
g_a=g(a),
\quad
g_m=g(m),
\quad
g_b=g(b),
\]
y
\[
q_a=q(a),
\quad
q_m=q(m),
\quad
q_b=q(b).
\]

El intervalo se subdivide si alguno de cuatro diagnosticos indica estructura
no resuelta.

\subsection{La derivada gira}

Se refina si
\begin{equation}
(g_m-g_a)(g_b-g_m)<0.
\label{eq:turning}
\end{equation}

\subsection{La curvatura cambia de signo}

Con
\[
\tau_q=10^{-8},
\]
se refina si
\begin{equation}
\min(q_a,q_m,q_b)<-\tau_q
\quad\text{y}\quad
\max(q_a,q_m,q_b)>\tau_q.
\label{eq:curvature-change}
\end{equation}

\subsection{La derivada central esta cerca de cero}

Sea
\[
\tau_g=10^{-8},
\qquad
\rho=0.2,
\]
y
\[
G_{\mathrm{edge}}
=
\max(|g_a|,|g_b|,\tau_g).
\]
Se refina si
\begin{equation}
|g_m|
\le
\max(
\tau_g,
\rho G_{\mathrm{edge}}
).
\label{eq:near-zero}
\end{equation}

\subsection{Cambio de signo oculto}

Tambien se refina si los extremos tienen el mismo signo pero una de las dos
mitades cambia de signo:
\[
\operatorname{sign}(g_a)
=
\operatorname{sign}(g_b),
\]
junto con
\[
g_ag_m<0
\quad\text{o}\quad
g_mg_b<0.
\]

\subsection{Parada}

El refinamiento se detiene al alcanzar
\[
\texttt{max\_depth}=8
\]
o cuando
\[
b-a\le10^{-3}.
\]

\section{Construccion de brackets}

Despues del refinamiento, los puntos evaluados se ordenan:
\[
S_1<S_2<\cdots<S_M.
\]

Un punto muestreado se considera raiz candidata si
\begin{equation}
|g(S_i)|\le10^{-8}.
\label{eq:near-root}
\end{equation}

Ademas, un par consecutivo define un bracket si
\begin{equation}
g(S_i)g(S_{i+1})<0.
\label{eq:bracket}
\end{equation}

Cada bracket se resuelve con Brent.

\section{Metodo de Brent}

La implementacion usa
\[
\texttt{scipy.optimize.brentq}
\]
para resolver
\[
g(S)=\widetilde F_S(S)=0.
\]

Brent conserva un bracket y combina tres mecanismos.

\subsection{Biseccion}

\[
S_{\mathrm{bis}}
=
\frac{a+b}{2}.
\]
Es el paso robusto que garantiza reducir el bracket.

\subsection{Secante}

Con dos puntos, la propuesta es
\begin{equation}
S_{\mathrm{sec}}
=
b
-
g(b)
\frac{b-a}{g(b)-g(a)}.
\label{eq:secant}
\end{equation}

\subsection{Interpolacion cuadratica inversa}

Con tres puntos adecuados, Brent aproxima \(S\) como funcion cuadratica de
\(g\) y evalua esa aproximacion en \(g=0\). Este paso puede ser mucho mas
rapido que biseccion cerca de una raiz regular.

\subsection{Salvaguarda}

Si una propuesta de secante o interpolacion no satisface las condiciones de
seguridad del metodo, se vuelve a biseccion. Asi se combina
\[
\boxed{
\text{robustez de bracket}
+
\text{velocidad de interpolacion}
}.
\]

La tolerancia actual es
\[
\texttt{root\_xtol}=10^{-10}.
\]

\section{Por que no Newton como buscador principal}

Newton sobre \(g(S)=0\) seria
\begin{equation}
S_{k+1}
=
S_k
-
\frac{g(S_k)}{q(S_k)}.
\label{eq:newton}
\end{equation}

Puede ser muy rapido cerca de una raiz, pero:

\begin{enumerate}[leftmargin=2.4em]
    \item necesita una semilla por cada raiz;
    \item puede saltar de una cuenca a otra;
    \item es problematico cuando \(q(S_k)\) es pequeno;
    \item no descubre por si solo cuantos puntos estacionarios existen.
\end{enumerate}

El objetivo de este proyecto es recuperar multiples minimos. Por eso se usa
\[
\text{refinamiento adaptativo}
\to
\text{brackets}
\to
\text{Brent}.
\]

\section{Deduplicacion y clasificacion}

Una misma raiz puede detectarse como punto casi estacionario y tambien mediante
un bracket. Despues de ordenar las raices se fusionan aquellas con separacion
menor que
\begin{equation}
\tau_{\mathrm{merge}}
=
\max(
10\,	exttt{root\_xtol},
10^{-8}
).
\label{eq:merge}
\end{equation}

Con los parametros actuales,
\[
\tau_{\mathrm{merge}}=10^{-8}.
\]

Para cada raiz deduplicada \(S^\star\), se evalua
\[
q(S^\star)
=
\widetilde F_{SS}(S^\star).
\]
Con \(\tau_q=10^{-8}\),
\begin{equation}
\begin{cases}
q(S^\star)>\tau_q
&
\Rightarrow
\texttt{minimum},
\\
q(S^\star)<-\tau_q
&
\Rightarrow
\texttt{maximum},
\\
|q(S^\star)|\le\tau_q
&
\Rightarrow
\texttt{flat}.
\end{cases}
\label{eq:classification}
\end{equation}

El tracking conserva los puntos \texttt{minimum} y los minimos de frontera.

\section{Minimos de frontera}

Los endpoints no necesitan satisfacer \(F_S=0\). El experimento evalua
\[
\widetilde F(0),
\qquad
\widetilde F(1),
\]
y usa
\[
\epsilon_b=0.002
\]
como sonda interior.

Se acepta \(S=0\) si
\begin{equation}
\widetilde F(0)
\le
\widetilde F(\epsilon_b),
\label{eq:left-boundary}
\end{equation}
y \(S=1\) si
\begin{equation}
\widetilde F(1)
\le
\widetilde F(1-\epsilon_b).
\label{eq:right-boundary}
\end{equation}

Estas son condiciones unilaterales de minimo.

\section{Separacion de candidatos}

Despues de reunir minimos interiores y de frontera, el experimento de tracking
elimina candidatos numericamente demasiado cercanos.

Se ordenan por menor Validation-1 loss. Un candidato \(S_b\) se descarta si ya
se retuvo \(S_a\) con
\begin{equation}
|S_b-S_a|
<
\epsilon_{\mathrm{cand}},
\qquad
\epsilon_{\mathrm{cand}}=0.02.
\label{eq:candidate-spacing}
\end{equation}

Esta operacion ocurre dentro de una sola superficie. No promedia smoothness ni
fusiona ramas temporales.

En el primer origin se inicializan como maximo cinco de esos minimos.

\section{Estabilidad numerica del espectro}

En aritmetica exacta,
\[
\operatorname{nullity}(D_d^\top D_d)=d.
\]
En punto flotante, los ceros estructurales pueden aparecer como numeros muy
pequenos positivos o negativos.

El codigo canonicaliza el espectro:
\[
\delta_1=\cdots=\delta_d=0,
\]
y para el resto
\[
\delta_j\leftarrow\max(\delta_j,0).
\]

Esto evita que un error de eigensolver en un valor que deberia ser cero sea
amplificado cuando \(\lambda\) es grande.

\section{Tratamiento de \(\lambda=\infty\)}

Para \(S=1\), no se calcula literalmente \(1+\infty\cdot0)^{-1}\. Se usa el
limite
\[
\alpha_j(\infty)
=
\begin{cases}
1,&\delta_j=0,\\
0,&\delta_j>0.
\end{cases}
\]
Las derivadas se fijan en cero. Asi
\[
\widehat{\bm\tau}_\infty
=
P_{\ker D_d}\bm x.
\]

\section{Cache de evaluaciones}

Durante la subdivision, un punto medio se convierte despues en extremo de dos
intervalos. Para evitar recalculo, se guarda
\[
S
\longmapsto
\left(
\lambda,
F,
F_\lambda,
F_{\lambda\lambda},
F_S,
F_{SS}
\right).
\]

Cada valor de \(S\) se evalua una sola vez. El contador
\texttt{n\_evaluations\_} mide precisamente cuantos smoothness distintos fueron
evaluados.

\section{Complejidad}

La eigendescomposicion densa inicial cuesta, en orden,
\[
O(L^3).
\]
Pero esa operacion se prepara una sola vez por configuracion.

Despues, cada nuevo \(\lambda\) reutiliza el espectro y el operador de forecast.
La evaluacion trabaja con pesos espectrales y multiplicaciones matriciales
vectorizadas sobre todos los rolling origins.

Por tanto, evaluar \(M\) candidatos no equivale a resolver desde cero
\(M\) problemas penalizados completos.

\section{Por que seguimos usando una grilla densa como diagnostico}

Ningun muestreo finito puede certificar, sin hipotesis adicionales, que
encontro todas las raices de una funcion suave arbitraria. Una oscilacion
suficientemente estrecha puede quedar entre puntos muestreados.

Por eso los benchmarks comparan
\[
\text{busqueda adaptativa derivativa}
\qquad\text{vs.}\qquad
\text{perfil denso de referencia}.
\]

La grilla densa sirve para revisar:

\begin{itemize}[leftmargin=2em]
    \item cuantos minimos se recuperaron;
    \item si se omitio una cuenca estrecha;
    \item el error de localizacion de cada \(S^\star\);
    \item el numero de evaluaciones ahorradas.
\end{itemize}

La grilla es una referencia numerica, no el algoritmo principal.

\section{Algoritmo completo}

Para una superficie fija:

\begin{enumerate}[leftmargin=2.4em]
    \item Preparar el objetivo rolling y su algebra espectral.
    \item Crear nueve puntos iniciales en \([0,1-10^{-6}]\).
    \item Agregar seis niveles diadicos cerca de ambos endpoints.
    \item Para cada nuevo \(S\):
    \begin{enumerate}
        \item invertir \(S\mapsto\lambda\) por biseccion;
        \item evaluar \(F,F_\lambda,F_{\lambda\lambda}\);
        \item transformar a \((F,F_S,F_{SS})\).
    \end{enumerate}
    \item Refinar los intervalos cuya derivada o curvatura sugiera estructura.
    \item Guardar puntos con \(|F_S|\le10^{-8}\).
    \item Crear brackets donde \(F_S\) cambia de signo.
    \item Resolver cada bracket con Brent hasta
    \(	exttt{xtol}=10^{-10}\).
    \item Deduplicar raices a escala \(10^{-8}\).
    \item Clasificar con \(F_{SS}\).
    \item Revisar \(S=0\) y \(S=1\) como minimos unilaterales.
    \item Separar candidatos dentro de \(0.02\).
    \item Inicializar hasta cinco ramas para el tracking temporal.
\end{enumerate}

\section{Flujo matematico}

El calculo de los minimos puede resumirse como

\[
\boxed{
S
\xrightarrow{\text{biseccion}}
\lambda(S)
\xrightarrow{\text{espectro}}
\widehat{\bm\tau},
\widehat{\bm\tau}',
\widehat{\bm\tau}''
}
\]

\[
\boxed{
\xrightarrow{C_{d,h,L}}
\widehat{\bm z},
\widehat{\bm z}',
\widehat{\bm z}''
\xrightarrow{\text{MSE}}
F,
F_\lambda,
F_{\lambda\lambda}
}
\]

\[
\boxed{
\xrightarrow{\text{regla de la cadena}}
F_S,
F_{SS}
\xrightarrow{\text{refinamiento}}
\text{brackets}
\xrightarrow{\text{Brent}}
S_j^\star
}.
\]

Finalmente,
\[
S_j^\star\in\mathcal M
\quad\Longleftrightarrow\quad
F_S(S_j^\star)=0
\quad\text{y}\quad
F_{SS}(S_j^\star)>0,
\]
salvo los minimos unilaterales de frontera.

\section{Relacion con el tracking temporal}

Todo lo anterior ocurre dentro de una sola ventana. El tracking temporal
documentado en
\texttt{tracked\_minimum\_smoothness\_selection.tex}
toma el conjunto
\[
\mathcal M_{d,r}
=
\{S_{d,1,r},\ldots,S_{d,J_r,r}\}
\]
y lo asocia con los minimos de la ventana siguiente mediante
\[
|S_{d,j,r+1}-S_{d,j,r}|
\le
\varepsilon_{\mathrm{track}}.
\]

Por tanto existen dos problemas distintos:

\[
\boxed{
\text{A: localizar todos los minimos de una superficie}
}
\]
y
\[
\boxed{
\text{B: asociar esos minimos entre ventanas}
}.
\]

Esta nota documenta A. La nota de tracking desarrolla B y la seleccion final
mediante Validation 2.

\section{Conclusion}

La seleccion numerica de smoothness no es una exploracion ciega de \(S\). La
implementacion explota la estructura lineal del suavizador para obtener
derivadas analiticas, transforma esas derivadas exactamente a la coordenada
normalizada, adapta el muestreo a la geometria observada y resuelve cada raiz
detectada con un metodo con bracket.

La idea central es
\[
\boxed{
\text{usar derivadas para descubrir donde buscar}
\quad+\quad
\text{usar Brent para localizar cada raiz robustamente}
}.
\]

Esto hace posible recuperar multiples minimos locales con muchas menos
evaluaciones que una grilla muy densa, mientras que los perfiles densos
permanecen como diagnostico independiente de cobertura.

\end{document}
